\documentclass[dvipdfmx,a4j]{jbook}
\usepackage{wrapfig}
\usepackage{tikz}
\usetikzlibrary{shapes,calc,arrows.meta}
\usepackage{amsmath}
\usepackage{amssymb}
\usepackage{fancyvrb}
\usepackage{tabularx}
%\graphicspath{{figs241004/}}
\usepackage{here}
\usepackage{graphicx}
\usepackage{fancyvrb}
\usepackage{multirow}
%\usepackage{amsthm}
\newtheorem{Def}{定義}[chapter]
\newtheorem{Thm}{定理}[chapter]
\newtheorem{Prop}{命題}[chapter]
\newtheorem{Exmp}{例}[chapter]
\newtheorem{Alg}{アルゴリズム}[chapter]
\newcommand{\tuple}[1]{\langle{#1}\rangle}
\newcommand{\Term}[1]{`\mbox{$\mathtt{#1}$}'}
\newcommand{\NTerm}[1]{\tuple{\mbox{$\mathsf{#1}$}}}
\newcommand{\First}[1]{\mathsf{First}_1({#1})}
\newcommand{\Follow}[1]{\mathsf{Follow}_1({#1})}
\newcommand{\greentext}[1]{{\color[rgb]{0,0.5,0}{#1}}}
\newcommand{\bluetext}[1]{{\color{blue}{#1}}}

\newcommand{\eof}{{\tt \#}}
\newcommand{\ttpl}{{\tt +}}
\newcommand{\ttsb}{{\tt -}}
\newcommand{\ttmu}{{\tt *}}
\newcommand{\ttdv}{{\tt /}}
\newcommand{\ttlp}{{\tt (}}
\newcommand{\ttrp}{{\tt )}}
\newcommand{\tteq}{{\tt =}}
\newcommand{\ttlt}{{\tt <}}
\newcommand{\ttcm}{{\tt ,}}

\newcommand{\ttif}{{\tt if}}
\newcommand{\ttfi}{{\tt fi}}
\newcommand{\ttthen}{{\tt then}}
\newcommand{\ttelse}{{\tt else}}
\newcommand{\ttwhile}{{\tt while}}
\newcommand{\ttdo}{{\tt do}}
\newcommand{\ttod}{{\tt od}}

\newcommand{\ttassgn}{{\tt :=}}
\newcommand{\ttsc}{{\tt ;}}

\newcommand{\ttand}{{\tt \&\&}}
\newcommand{\ttor}{{\tt ||}}
\newcommand{\ttneg}{{\tt !}}

\newcommand{\rppush}[1]{\mathsf{push}\ \texttt{#1}}
\newcommand{\rpassign}{{\tt :=}}
\newcommand{\rpframe}[1]{\draw #1 -- ($#1+(0,1)$) #1 -- ($#1+(1.5,0)$) ($#1+(1.5,0)$) -- ($#1+(1.5,1)$);}
\newcommand{\rpframerow}[2]{\draw ($#1}

\newcommand{\newTemp}[1]{{\tt newTemp(}#1\texttt{)}}
\newcommand{\newTempnull}{{\tt newTemp}}
\newcommand{\newLabel}[1]{{\tt newLabel(}#1\texttt{)}}
\newcommand{\newLabelnull}{{\tt newLabel}}
\newcommand{\generatefo}[4]{{\tt generate(}{#1},=,{#2},{#3},{#4}\texttt{)}}
\newcommand{\generateth}[3]{{\tt generate(}{#1},=,{#2},{#3}\texttt{)}}
\newcommand{\generatetw}[2]{{\tt generate(}{#1},=,{#2}\texttt{)}}
\newcommand{\generateon}[1]{{\tt generate(}{#1}\texttt{)}}
\newcommand{\generateCTL}[3]{{\tt generate(}{#1}.{#2},{#3}\texttt{)}}
\newcommand{\generate}[1]{{\tt generate(}{#1}{\tt )}}
\newcommand{\generatenull}{{\tt generate}}
\newcommand{\generateassgin}[2]{{\tt generate}(`{#1}{\tt =}{#2}')}
\newcommand{\generateif}[2]{{\tt generate}(`{\tt if}\ {#1}\ {\tt goto}\ {#2}')}
\newcommand{\generateifnot}[2]{{\tt generate}(`{\tt if\_not}\ {#1}\ {\tt goto}\ {#2}')}
\newcommand{\generatelabel}[1]{{\tt generate(`label\ }{#1}')}
\newcommand{\generategoto}[1]{{\tt generate(`goto\ }{#1}')}

\newcolumntype{b}{X}
\newcolumntype{s}{>{\hsize=.5\hsize}X}

\usepackage{subfig}
%\usepackage[framemethod=tikz]{mdframed}
\newcounter{coffee}
\newcounter{bean}
\usepackage{docmute}
\usepackage{url}
\usepackage{makeidx}
\usepackage{here}
\usepackage{array}
\usepackage{tabularx}
\usepackage{mdframed}
\usepackage{longtable}
\newtheorem{Def}{定義}[chapter]
\newtheorem{Thm}{定理}[chapter]
\newtheorem{Prop}{命題}[chapter]
\newtheorem{Lem}{補題}[chapter]
\newtheorem{Exmp}{例}[chapter]
\newtheorem{Alg}{アルゴリズム}[chapter]

\newtheorem*{repeatedtheorem}{定理 \repeatednumber}
\newtheorem*{repeatedlemma}{補題 \repeatednumber}

\renewcommand{\proofname}{証明}

\newenvironment{repThm}[1]
{\def\repeatednumber{\ref{#1}}\begin{repeatedtheorem}}
{\end{repeatedtheorem}}

\newenvironment{repLem}[1]
{\def\repeatednumber{\ref{#1}}\begin{repeatedlemma}}
{\end{repeatedlemma}}

\newcommand{\tuple}[1]{\langle{#1}\rangle}
\newcommand{\Term}[1]{`\mbox{$\mathtt{#1}$}'}
\newcommand{\NTerm}[1]{\tuple{\mbox{$\mathsf{#1}$}}}
\newcommand{\First}[1]{\mathsf{First}_1({#1})}
\newcommand{\Follow}[1]{\mathsf{Follow}_1({#1})}
\newcommand{\greentext}[1]{{\color[rgb]{0,0.5,0}{#1}}}
\newcommand{\bluetext}[1]{{\color{blue}{#1}}}

\newcommand{\eof}{{\tt \#}}
\newcommand{\ttpl}{{\tt +}}
\newcommand{\ttsb}{{\tt -}}
\newcommand{\ttmu}{{\tt *}}
\newcommand{\ttdv}{{\tt /}}
\newcommand{\ttlp}{{\tt (}}
\newcommand{\ttrp}{{\tt )}}
\newcommand{\tteq}{{\tt =}}
\newcommand{\ttlt}{{\tt <}}
\newcommand{\ttcm}{{\tt ,}}

\newcommand{\ttbrkl}{{\tt [}}
\newcommand{\ttbrkr}{{\tt ]}}

\newcommand{\ttif}{{\tt if}}
\newcommand{\ttfi}{{\tt fi}}
\newcommand{\ttthen}{{\tt then}}
\newcommand{\ttelse}{{\tt else}}
\newcommand{\ttwhile}{{\tt while}}
\newcommand{\ttdo}{{\tt do}}
\newcommand{\ttod}{{\tt od}}

\newcommand{\ttassgn}{{\tt :=}}
\newcommand{\ttsc}{{\tt ;}}

\newcommand{\ttand}{{\tt \&\&}}
\newcommand{\ttor}{{\tt ||}}
\newcommand{\ttneg}{{\tt !}}

\newcommand{\tti}{\texttt{i}}
\newcommand{\ttj}{\texttt{j}}
\newcommand{\ttk}{\texttt{k}}
\newcommand{\ttx}{\texttt{x}}
\newcommand{\tty}{\texttt{y}}
\newcommand{\ttz}{\texttt{z}}

\newcommand{\rppush}[1]{\mathsf{push}\ \texttt{#1}}
\newcommand{\rpassign}{{\tt :=}}
\newcommand{\rpframe}[1]{\draw #1 -- ($#1+(0,1)$) #1 -- ($#1+(1.5,0)$) ($#1+(1.5,0)$) -- ($#1+(1.5,1)$);}
\newcommand{\rpframerow}[2]{\draw ($#1}

\newcommand{\newTemp}[1]{{\tt newTemp(}#1\texttt{)}}
\newcommand{\newTempnull}{{\tt newTemp}}
\newcommand{\newLabel}[1]{{\tt newLabel(}#1\texttt{)}}
\newcommand{\newLabelnull}{{\tt newLabel}}
\newcommand{\generatefo}[4]{{\tt generate(}{#1},=,{#2},{#3},{#4}\texttt{)}}
\newcommand{\generateth}[3]{{\tt generate(}{#1},=,{#2},{#3}\texttt{)}}
\newcommand{\generatetw}[2]{{\tt generate(}{#1},=,{#2}\texttt{)}}
\newcommand{\generateon}[1]{{\tt generate(}{#1}\texttt{)}}
\newcommand{\generateCTL}[3]{{\tt generate(}{#1}.{#2},{#3}\texttt{)}}
\newcommand{\generate}[1]{{\tt generate(}{#1}{\tt )}}
\newcommand{\generatenull}{{\tt generate}}
\newcommand{\generateassgin}[2]{{\tt generate}(`{#1}{\tt =}{#2}')}
\newcommand{\generateif}[2]{{\tt generate}(`{\tt if}\ {#1}\ {\tt goto}\ {#2}')}
\newcommand{\generateifnot}[2]{{\tt generate}(`{\tt if\_not}\ {#1}\ {\tt goto}\ {#2}')}
\newcommand{\generatelabel}[1]{{\tt generate(`label\ }{#1}')}
\newcommand{\generategoto}[1]{{\tt generate(`goto\ }{#1}')}

\newcolumntype{b}{X}
\newcolumntype{s}{>{\hsize=.5\hsize}X}

\newcolumntype{Y}{>{\centering\arraybackslash}X}
\newcommand{\inst}[1]{\begin{tabularx}{.3\textwidth}{|Y|Y|Y|Y|}\hline\texttt{#1}& $b$ & $i$ &$a$\\\hline\end{tabularx}}
\newcommand{\ignored}[1]{\par #1は無視される。}

\newcommand{\bnfnts}[1]{\langle\mathrm{#1}\rangle}
\newcommand{\bnfts}[1]{\texttt{'#1'}}
\newcommand{\bnfarr}{\texttt{::=}}

\newcommand{\domv}[1]{\mathsf{dom}({#1})}
\newcommand{\idomv}[1]{\mathsf{idom}({#1})}
\newcommand{\domgg}{\underline{\gg}}
\newcommand{\domll}{\underline{\ll}}

\newcommand{\loadinst}[2]{\texttt{Load}\>{$#1$}\>{$#2$}}
\newcommand{\oprinst}[3]{\texttt{#1}\>{#2}\>{#3}}

\setcounter{chapter}{0}
\begin{document}
\appendix
\chapter{VMコードの仕様}

6章にあげたスタックマシンの目的コードを拡張して，手続き呼び出しを可能にする．
このために10$\sim$12の３つの命令を追加する．

命令長は1 wordであり、命令は次の形式を持つ。
\begin{center}
\begin{tabularx}{.3\textwidth}{|c|Y|Y|Y|Y|}\hline
    OPCODE & $b$ & $i$ & $a$\\\hline
\end{tabularx}
\end{center}
% \[
% \begin{array}{|c|c|c|c|}\hline
% \mathrm{OPCODE}&b&i&a\\\hline
% \end{array}
% \]
\begin{description}
\item[$b$] ベースレジスタの指定。
\item[$i$] インデックスレジスタの指定。
\item[$a$] アドレス（displacement）。命令により定数や演算の指定にも使う。
\end{description}
命令領域（読出し専用）はword単位で番地を持ち、0番地から実行する。
PCが示す番地の命令をIRに格納し、PCを1増加させてから、IRの命令を実行する。
初期状態は次のとおりである。
\[
\mathrm{PC}=0,\quad R_1=0,\quad R_2=-1,\quad R_3=0,\quad R_4=0.
\]

\paragraph{1. Load literal}
\inst{lit}
\begin{verbatim}
R2 = R2 + 1; s[R2] = a
\end{verbatim}
\ignored{$b$と$i$}
\noindent 原文の\texttt{'a'}は即値$a$を表す。

\paragraph{2. Load}
\inst{lod}
\begin{verbatim}
R2 = R2 + 1; s[R2] = s[Rb + Ri + a]
\end{verbatim}
$b$と$i$が指定するレジスタの内容と$a$を加えたものをオペランドアドレスとする。
ただし、$b=0$のときベースによるアドレス修飾は行わず、$i=0$のときも同様である。

\paragraph{3. Store}
\inst{sto}
\begin{verbatim}
s[Rb + Ri + a] = s[R2]; R2 = R2 - 1
\end{verbatim}
アドレスの計算は\texttt{lod}と同様であり、$b=0$または$i=0$の項は加えない。

\paragraph{4. Move to index register}
\inst{mvx}
\begin{verbatim}
Ri = s[R2]; R2 = R2 - 1
\end{verbatim}
\ignored{$b$と$a$}

\paragraph{5. Arithmetic operation}
\inst{opr}
\par $a$により次の11種類の演算を指定する。
\ignored{$b$と$i$}
\begin{longtable}{|c|p{0.83\linewidth}|}
\hline $a$ & 動作\\\hline
\endhead
0 & \texttt{s[R2] = -s[R2]}\\\hline
1 & \texttt{R2 = R2 - 1; s[R2] = s[R2] + s[R2+1]}\\\hline
2 & \texttt{R2 = R2 - 1; s[R2] = s[R2] - s[R2+1]}\\\hline
3 & \texttt{R2 = R2 - 1; s[R2] = s[R2] * s[R2+1]}\\\hline
4 & \texttt{R2 = R2 - 1; s[R2] = s[R2] div s[R2+1]}\\\hline
5 & \texttt{R2 = R2 - 1;} $s[R_2]=s[R_2+1]$なら\texttt{s[R2] = 1}、そうでなければ\texttt{s[R2] = 0}。\\\hline
6 & \texttt{R2 = R2 - 1;} $s[R_2]\ne s[R_2+1]$なら\texttt{s[R2] = 1}、そうでなければ\texttt{s[R2] = 0}。\\\hline
7 & \texttt{R2 = R2 - 1;} $s[R_2]<s[R_2+1]$なら\texttt{s[R2] = 1}、そうでなければ\texttt{s[R2] = 0}。\\\hline
8 & \texttt{R2 = R2 - 1;} $s[R_2]\le s[R_2+1]$なら\texttt{s[R2] = 1}、そうでなければ\texttt{s[R2] = 0}。\\\hline
9 & \texttt{R2 = R2 - 1;} $s[R_2]>s[R_2+1]$なら\texttt{s[R2] = 1}、そうでなければ\texttt{s[R2] = 0}。\\\hline
10 & \texttt{R2 = R2 - 1;} $s[R_2]\ge s[R_2+1]$なら\texttt{s[R2] = 1}、そうでなければ\texttt{s[R2] = 0}。\\\hline
\end{longtable}

\paragraph{6. Unconditional Jump}
\inst{jmp}
\begin{verbatim}
PC = a
\end{verbatim}
\ignored{$b$と$i$}

\paragraph{7. Jump on condition}
\inst{jpc}
\begin{verbatim}
if s[R2] = 0 then PC = a;
R2 = R2 - 1
\end{verbatim}
\ignored{$b$と$i$}

\paragraph{8. Read data}
\inst{get}
\par \texttt{R2 = R2 + 1}として、標準入力から次のデータ（整数値）を読み、
\texttt{s[R2]}に格納する。
\ignored{$b$、$i$、$a$}

\paragraph{9. Write data}
\inst{put}
\par \texttt{s[R2]}の内容を標準出力に出力し、\texttt{R2 = R2 - 1}とする。
\ignored{$b$、$i$、$a$}

\paragraph{10. Increment stack top}
\inst{int}
\begin{verbatim}
R2 = R2 + a
\end{verbatim}
\ignored{$b$と$i$}

\paragraph{11. Procedure call}
\inst{cal}
\begin{verbatim}
s[R2+1] = R1; s[R2+2] = PC;
s[R2+3] = R3; s[R2+4] = R4;
R1 = R2 + 5; R2 = R2 + 4; PC = a
\end{verbatim}
\ignored{$b$と$i$}

\paragraph{12. Return}
\inst{rtn}
\begin{verbatim}
R2 = R1 - 5; PC = s[R1-3];
R3 = s[R1-2]; R4 = s[R1-1]; R1 = s[R1-4]
\end{verbatim}
\ignored{$b$、$i$、$a$}

\end{document}
